\chapter*{TÓM TẮT}\addcontentsline{toc}{chapter}{Tóm tắt, Abstract và Từ khóa}
Nghiên cứu nhằm nhận diện và đo lường các yếu tố ảnh hưởng đến quyết định chọn trường và ngành học của học sinh THPT tại các tỉnh Đồng bằng sông Cửu Long. Trên nền tảng mô hình lựa chọn trường của Chapman, lý thuyết hành vi có kế hoạch và các nghiên cứu tại Việt Nam, mô hình đề xuất gồm bảy nhóm yếu tố: phù hợp cá nhân; triển vọng nghề nghiệp; danh tiếng và chất lượng đào tạo; chi phí và hỗ trợ tài chính; ảnh hưởng xã hội; truyền thông tuyển sinh; vị trí và điều kiện học tập. Thiết kế nghiên cứu kết hợp định tính hiệu chỉnh thang đo và định lượng cắt ngang. Kế hoạch lấy mẫu phân tầng theo địa bàn và loại trường, với cỡ mẫu tối thiểu 385 và mục tiêu 500 học sinh lớp 12. Dữ liệu được kiểm định bằng Cronbach's Alpha, EFA, tương quan và hồi quy đa biến.

Kết quả trên 486 quan sát cho thấy phù hợp cá nhân và triển vọng nghề nghiệp có tác động mạnh nhất, tiếp theo là danh tiếng--chất lượng, chi phí, ảnh hưởng xã hội, vị trí và truyền thông. Từ đó, nghiên cứu đề xuất các hàm ý cho cơ sở giáo dục đại học, trường THPT, gia đình và học sinh trong hoạt động tư vấn hướng nghiệp và tuyển sinh có căn cứ.

\noindent\textbf{Từ khóa:} chọn trường đại học, chọn ngành học, học sinh THPT, hành vi lựa chọn, Đồng bằng sông Cửu Long.

\chapter*{ABSTRACT}
This study identifies and measures factors influencing high-school students' university and major choice in the Vietnamese Mekong Delta. Drawing on Chapman's college-choice model, the Theory of Planned Behavior, and Vietnamese evidence, the proposed model contains seven constructs: person--major fit, career prospects, institutional reputation and quality, cost and financial support, social influence, admission communication, and location and learning conditions. A sequential design combines qualitative scale refinement with a cross-sectional survey. Stratified sampling by province and school type targets 500 Grade-12 students, with a minimum sample of 385. Reliability analysis, exploratory factor analysis, correlation, and multiple regression are proposed.

Results from 486 observations show that person--major fit and career prospects exert the strongest effects, followed by reputation--quality, costs, social influence, location, and communication. Practical implications are offered for universities, high schools, families, and students.

\noindent\textbf{Keywords:} university choice, major choice, high-school students, choice behavior, Mekong Delta.
